\section*{الفصل \ref{ch:g9:acids-bases-ph} --- الأحماض والقواعد وpH}

\begin{solution}{exo:g9:acids-bases-ph:1}
\omterm{def:g9:acids-bases-ph:ph}{pH} 3: حمضي. \omterm{def:g9:acids-bases-ph:ph}{pH} 7: متعادل. \omterm{def:g9:acids-bases-ph:ph}{pH} 11: قاعدي.
\end{solution}

\begin{solution}{exo:g9:acids-bases-ph:2}
الحمضي: \omterm{def:g9:ions:ion}{أيونات} \ce{H+}. القاعدي: \omterm{def:g9:ions:ion}{أيونات} \ce{OH-}.
\end{solution}

\begin{solution}{exo:g9:acids-bases-ph:3}
عصير الليمون (2)، والقهوة (5)، والماء النقي (7)، وماء البحر (8.1)، والنشادر
المنزلي (12).
\end{solution}

\begin{solution}{exo:g9:acids-bases-ph:4}
يوضع شريط في صحن نظيف جاف، ويُلمس بقطرة من \omterm{def:g2:mixing-and-dissolving:dissolve}{المحلول} تحملها ساق زجاجية
نظيفة؛ ويُقارن اللون فورًا بسلّم
الألوان.
\end{solution}

\begin{solution}{exo:g9:acids-bases-ph:5}
نحو 5 بعد تخفيف بعشرة أضعاف؛ ونحو 6 بعد تخفيف بمئة ضعف.
\end{solution}

\begin{solution}{exo:g9:acids-bases-ph:6}
لا. فالتخفيف يقرّب \omterm{def:g9:acids-bases-ph:ph}{pH} من 7 لكنه لا يتجاوزه أبدًا: فالماء نفسه متعادل،
وإضافته لا يمكن أن تجعل \omterm{def:g9:ions:ion}{أيونات} \ce{OH-} أكثر عددًا من \omterm{def:g9:ions:ion}{أيونات} \ce{H+}.
\end{solution}

\begin{solution}{exo:g9:acids-bases-ph:7}
\omterm{def:g2:mixing-and-dissolving:mix}{الخلط} يطلق حرارة. فالماء الذي يُصب على حمض مركّز قد يغلي في الحال ويقذف
الحمض؛ أما الحمض الذي يُصب ببطء في الماء فينتشر فورًا في حجم
كبير.
\end{solution}

\begin{solution}{exo:g9:acids-bases-ph:8}
أكّال: يحرق الجلد والعينين، ويهاجم الفلزات. البس نظارات واقية وقفازات؛
واشطف أي رذاذ بكثير من الماء.
\end{solution}

\begin{solution}{exo:g9:acids-bases-ph:9}
ثلاثة تخفيفات (من 1 إلى 2، ومن 2 إلى 3، ومن 3 إلى 4): $10 \times 10 \times 10 =
1000$ مرة.
\end{solution}

\begin{solution}{exo:g9:acids-bases-ph:10}
حليب المغنيسيا قاعدي (\omterm{def:g9:acids-bases-ph:ph}{pH} له نحو 10.5): فهو يتفاعل مع جزء من الحمض الزائد
في المعدة.
\end{solution}

\begin{solution}{exo:g9:acids-bases-ph:11}
لا: فقيمة \omterm{def:g9:acids-bases-ph:ph}{pH} له 8.1 تعني أنه لا يزال قاعديًا قليلًا. وهو يسير نحو \omterm{def:g9:acids-bases-ph:ph}{pH} أخفض،
أي أقل قاعدية، لأن ثنائي أكسيد الكربون من \omterm{def:g4:air-a-mixture-of-gases:air}{الهواء} \omterm{def:g2:mixing-and-dissolving:dissolve}{يذوب}
فيه.
\end{solution}

\begin{solution}{exo:g9:acids-bases-ph:12}
\omterm{def:g9:acids-bases-ph:ph-paper}{ورق pH} يُقرأ بالعين مقابل سلّم ألوان، بدقة وحدة واحدة تقريبًا؛ أما \omterm{def:g9:acids-bases-ph:ph-paper}{مقياس pH}
فيعطي عددًا بمنزلة عشرية أو منزلتين. فالمقياس أدق؛ والنتيجتان متفقتان في
حدود دقة الورق.
\end{solution}

\begin{solution}{pb:g9:acids-bases-ph:1}
\textbf{1.} حمضي: يحوي في الغالب \omterm{def:g9:ions:ion}{أيونات} \ce{H+(aq)} و\ce{Cl-(aq)}.

\textbf{2.} GHS05، أكّال: فالحمض يهاجم الجلد والعينين والفلزات.

\textbf{3.} يعطي عددًا دقيقًا بدلًا من لون يُقارن
بالعين.

\textbf{4.} عشر مرات: فوحدة \omterm{def:g9:acids-bases-ph:ph}{pH} واحدة توافق تخفيفًا بعشرة أضعاف.

\textbf{5.} \qty{10}{L}، وقيمة \omterm{def:g9:acids-bases-ph:ph}{pH} له 3.

\textbf{6.} ثلاثة: من \omterm{def:g9:acids-bases-ph:ph}{pH} 2 إلى 3، ومن 3 إلى 4، ومن 4 إلى 5.

\textbf{7.} $1 \times 10 \times 10 \times 10 = \qty{1000}{L}$.

\textbf{8.} لا: فالتخفيف لا يفعل إلا أن يقرّب \omterm{def:g9:acids-bases-ph:ph}{pH} من 7.

\textbf{9.} \ce{H+ + OH- -> H2O}.

\textbf{10.} يزول الحمض، إذ يتحول ماءً، بدلًا من أن ينتشر في حجم هائل؛ ولا
يلزم إلا قليل من القاعدة، لا ألف لتر من الماء.

\textbf{11.} صودا الخبز المذابة في الماء (أو الماء والصابون).

\textbf{12.} $1000 - 1 = \qty{999}{L}$ من الماء.
\end{solution}
